\section{Combination Sum}
\label{sec:rec-combination-sum}

\problemstatement{Given an array $\textit{candidates}$ of \textbf{distinct}
positive integers and a target $\textit{target}$, return every unique
combination of candidates that sums to $\textit{target}$. The
\textbf{same candidate may be used an unlimited number of times} within
one combination.}

\begin{patternbox}
\emph{``The same element may be reused unlimited times''} is the
keyword that distinguishes this from Section~\ref{sec:rec-power-set}'s
plain pick/not-pick tree: the \textbf{pick} branch must \textbf{stay at
the same index} (since this element remains available for reuse) rather
than advancing to the next index, while the \textbf{not-pick} branch
still advances (moving permanently past this candidate, never to
reconsider it).
\end{patternbox}

\subsection*{Understanding the problem carefully}

At each index, we still face a binary choice -- use this candidate (and
possibly use it again later) or move past it forever -- but ``use this
candidate'' no longer means ``use it exactly once and move on'': because
reuse is unlimited, choosing to include $\textit{candidates}[\textit{index}]$
should recurse back into the \emph{same} index, leaving open the
possibility of including it again. The recursion only definitively moves
past an index once the not-pick branch is taken.

\begin{ideabox}[Pick Stays, Not-Pick Advances]
$\textsc{Solve}(\textit{index}, \textit{remaining}, \textit{current})$:
if $\textit{remaining}=0$: record $\textit{current}$; return. If
$\textit{index}=n$ or $\textit{remaining}<0$: return (failed branch,
either ran out of candidates or overshot the target). Otherwise: append
$\textit{candidates}[\textit{index}]$ to $\textit{current}$; recurse
$\textsc{Solve}(\textit{index}, \textit{remaining} -
\textit{candidates}[\textit{index}], \textit{current})$ (pick, staying
at the same index); remove the appended element (undo); recurse
$\textsc{Solve}(\textit{index}+1, \textit{remaining}, \textit{current})$
(not-pick, advancing).
\end{ideabox}

\subsection*{Why staying at the same index correctly models unlimited reuse}

If the pick branch advanced to $\textit{index}+1$ as in
Section~\ref{sec:rec-power-set}, each candidate could be used at most
once per combination -- exactly the wrong behavior for this problem.
Staying at $\textit{index}$ after picking means the very next recursive
call still has $\textit{candidates}[\textit{index}]$ available to pick
\emph{again}, and can keep doing so for as many repetitions as the
remaining target allows, only ever moving past index $\textit{index}$
once the not-pick branch is explicitly taken (which represents
``permanently done with this candidate, for this branch of the search'').
The $\textit{remaining}<0$ check ensures the recursion does not loop
forever picking a positive candidate repeatedly past the point where it
could still sum to the target -- it is the base-case escape that makes
termination certain despite the same index being revisited arbitrarily
many times.

\subsection*{Algorithm}

\begin{enumerate}
  \item $\textsc{Solve}(\textit{index}, \textit{remaining}, \textit{current})$:
  \begin{enumerate}
    \item If $\textit{remaining}=0$: record $\textit{current}$; return.
    \item If $\textit{index}=n$ or $\textit{remaining}<0$: return.
    \item Append $\textit{candidates}[\textit{index}]$;
          $\textsc{Solve}(\textit{index}, \textit{remaining} -
          \textit{candidates}[\textit{index}], \textit{current})$;
          remove (undo).
    \item $\textsc{Solve}(\textit{index}+1, \textit{remaining},
          \textit{current})$.
  \end{enumerate}
  \item Call $\textsc{Solve}(0, \textit{target}, [\,])$.
\end{enumerate}

\subsection*{Dry run}

\dryrun{Take $\textit{candidates} = [2,3]$, $\textit{target}=5$.}

\begin{itemize}
  \item $\textsc{Solve}(0,5,[])$: pick $2$:
        $\textsc{Solve}(0,3,[2])$ (stays at index $0$).
  \begin{itemize}
    \item Pick $2$ again: $\textsc{Solve}(0,1,[2,2])$. Pick $2$ again:
          $\textsc{Solve}(0,-1,[2,2,2])$: $\textit{remaining}<0$,
          return (fail). Undo. Not-pick: $\textsc{Solve}(1,1,[2,2])$:
          pick $3$: $\textsc{Solve}(1,-2,\ldots)$: fail. Undo.
          Not-pick: $\textsc{Solve}(2,1,[2,2])$: $\textit{index}=n=2$,
          return (fail).
    \item Undo back to $[2]$. Not-pick $3$ at this level:
          $\textsc{Solve}(1,3,[2])$: pick $3$:
          $\textsc{Solve}(1,0,[2,3])$: $\textit{remaining}=0$! Record
          $[2,3]$.
  \end{itemize}
  \item Back at top: undo to $[]$. Not-pick $2$:
        $\textsc{Solve}(1,5,[])$: pick $3$: $\textsc{Solve}(1,2,[3])$:
        pick $3$ again: $\textsc{Solve}(1,-1,[3,3])$: fail. Not-pick:
        $\textsc{Solve}(2,2,[3])$: $\textit{index}=n$, fail.
\end{itemize}

The only recorded combination is $[2,3]$. (Note: in a full trace,
$[2,2,\ldots]$ paths never reach exactly $0$ remaining and so never
record anything, which the dry run above confirms.)

\subsection*{C++ implementation}

\begin{lstlisting}
#include <bits/stdc++.h>
using namespace std;

void solve(int index, int remaining, vector<int>& candidates,
           vector<int>& current, vector<vector<int>>& result) {
    if (remaining == 0) {
        result.push_back(current);
        return;
    }
    if (index == (int)candidates.size() || remaining < 0) return;

    current.push_back(candidates[index]);
    solve(index, remaining - candidates[index], candidates, current, result);
    current.pop_back();

    solve(index + 1, remaining, candidates, current, result);
}

vector<vector<int>> combinationSum(vector<int>& candidates, int target) {
    vector<vector<int>> result;
    vector<int> current;
    solve(0, target, candidates, current, result);
    return result;
}

int main() {
    vector<int> candidates = {2, 3, 6, 7};
    for (auto& comb : combinationSum(candidates, 7)) {
        cout << "[ ";
        for (int x : comb) cout << x << " ";
        cout << "] ";
    }
    cout << endl;
    return 0;
}
\end{lstlisting}

\begin{complexitybox}
\textbf{Time Complexity.} Difficult to bound tightly in general (it
depends on the target and candidate values), but is exponential in the
worst case, roughly $O(2^{\textit{target}})$ for the branching factor at
each remaining-target unit; in practice it is bounded by the number of
valid combinations times their average length.
\textbf{Space Complexity.} $O(\textit{target}/\min(\textit{candidates}))$
for the maximum recursion depth (the most times the smallest candidate
could be picked before exceeding the target).
\end{complexitybox}

\begin{takeawaybox}
``Unlimited reuse of the same element'' converts the pick branch's
target index from ``advance to $\textit{index}+1$'' to ``stay at
$\textit{index}$.'' This single change is the entire difference between
Combination Sum (this section) and the no-reuse subset/combination
problems before and after it -- worth recognizing on sight whenever a
problem statement explicitly allows repeated use of the same value.
\end{takeawaybox}
